% Preamble
\documentclass[11pt, aspectratio=169]{beamer}
	% Packages
	\usepackage[english]{babel}
	\usepackage{lipsum}
	\usepackage{mathptmx}
	\usepackage{xcolor}
	% Themes
	\usefonttheme[onlymath]{serif}
	\usetheme[nofirafonts]{focus}
	\renewcommand{\thempfootnote}{\arabic{mpfootnote}}
	% Cover
	\title{Microeconometrics Using Stata}
	\subtitle{BINARY OUTCOME MODELS}
	\author{
		Jhon R. Ordoñez
		\footnote{\label{1} National University of San Cristóbal de Huamanga}
		\footnote{\label{2} Faculty of Economic, Administrative and Accounting Sciences}
		\footnote{\label{3} Professional School of Economics}
	}
	\date{\today}
% Body
\begin{document}
	% Cover
	\begin{frame}
		\maketitle
	\end{frame}
	% Outline
	\begin{frame}[t]
		\frametitle{Outline}
		\tableofcontents
	\end{frame}
	% Section 1
	\section{Exercise 1}
		% Slide 1
		\begin{frame}
			\frametitle{Exercise 1}
				\begin{block}{Exercise 1}
					Consider the example of \textcolor{red!50}{\textbf{section 17.4}} with dependent variable \textcolor{blue}{\textbf{ins}} and
					the single regressor \textcolor{blue}{\textbf{educyear}}. Estimate the parameters of logit, probit,
					and \textbf{OLS} models using both default and robust standard errors. For the
					regressor \textcolor{blue}{\textbf{educyear}}, compare its coefficient across the models,
					compare default and robust standard errors of this coefficient, and
					compare the $t$ statistics based on robust standard errors. For each
					model, compute the \textbf{ME} of one more year of education for someone
					with sample mean years of education, as well as the \textbf{AME}. Which
					model fits the data better-logit or probit?	
				\end{block}
		\end{frame}
	% Section 2
	\section{Exercise 2}
		% Slide 1
		\begin{frame}
			\frametitle{Exercise 2}
			\begin{block}{Exercise 2}
				Use the \textcolor{blue}{\textbf{cloglog}} command to estimate the parameters of the binary
				probability model for \textcolor{blue}{\textbf{ins}} with the same explanatory variables used in
				the \textcolor{blue}{\textbf{logit}} model in this chapter. Estimate the AME for the regressors.
				Calculate the odds ratios of \textcolor{blue}{\textbf{ins=1}} for the following values of the
				covariates: \textcolor{blue}{\textbf{age=50}}, \textcolor{blue}{\textbf{retire=0}}, \textcolor{blue}{\textbf{hstatusg=1}}, \textcolor{blue}{\textbf{hhincome=45}},
				\textcolor{blue}{\textbf{educyear=12}}, \textcolor{blue}{\textbf{married=1}}, and \textcolor{blue}{\textbf{hisp=0}}.
			\end{block}
		\end{frame}
	% Section 3
	\section{Exercise 3}
		% Slide 1
		\begin{frame}
			\frametitle{Exercise 3}
			\begin{block}{Exercise 3}
				Generate a graph of fitted probabilities against years of education
				(\textcolor{blue}{\textbf{educyear}}) or age (\textcolor{blue}{\textbf{age}}) using as a template the commands used for
				generating \textcolor{green!75}{\textbf{figure 17.1}} in this chapter.
			\end{block}
		\end{frame}
	% Section 4
	\section{Exercise 4}
		% Slide 1
		\begin{frame}
			\frametitle{Exercise 4}
			\begin{block}{Exercise 4}
				Estimate the parameters of the logit model of \textcolor{red!50}{\textbf{section 17.4.2.}} Now,
				estimate the parameters of the probit model using the \textcolor{blue}{\textbf{probit}}
				command. Use the reported log likelihoods to compare the models by
				the Akaike’s information criterion and \textbf{BIC}.
			\end{block}
		\end{frame}
	% Section 5
	\section{Exercise 5}
		% Slide 1
		\begin{frame}
			\frametitle{Exercise 5}
			\begin{block}{Exercise 5}
				Estimate the probit regression of \textcolor{red!50}{\textbf{section 17.4.4.}} Using the
				conditioning values (\textcolor{blue}{\textbf{age=65}}, \textcolor{blue}{\textbf{retire=1}}, \textcolor{blue}{\textbf{hstatusg=1}}, \textcolor{blue}{\textbf{hhincome=60}},
				\textcolor{blue}{educyear=17}, \textcolor{blue}{\textbf{married=1}}, \textcolor{blue}{\textbf{hisp=0}}), estimate and compare the \textbf{ME} of
				\textcolor{blue}{\textbf{age}} on the Pr(\textcolor{blue}{\textbf{ins=1|x}}), using both the \textcolor{blue}{\textbf{margins}} and \textcolor{blue}{\textbf{prchange}}
				commands. They should give the same result.
			\end{block}
		\end{frame}
	% Section 6
	\section{Exercise 6}
		% Slide 1
		\begin{frame}
			\frametitle{Exercise 6}
			\begin{block}{Exercise 6}
				Using the \textcolor{blue}{\textbf{hetprobit}} command, estimate the parameters of the model
				of \textcolor{red!50}{\textbf{section 17.4}}, using \textcolor{blue}{\textbf{hhincome}} as the sole variable determining the
				variance. Test the null hypothesis of homoskedastic probit.
			\end{block}
		\end{frame}
	% Section 7
	\section{Exercise 7}
		% Slide 1
		\begin{frame}
			\frametitle{Exercise 7}
			\begin{block}{Exercise 7}
				Using the example in \textcolor{red!50}{\textbf{section 17.10}} as a template, estimate a grouped
				logistic regression using \textcolor{blue}{\textbf{educyear}} as the grouping variable. Comment
				on what you regard as unsatisfactory features of the grouping variable
				and the results.
			\end{block}
		\end{frame}
	% References
	\begin{frame}[t,allowframebreaks]
		\frametitle{References}
		\bibliography{biblio.bib}
		% Style
		\bibliographystyle{apalike}
		% Authors
		\nocite{cameron2022-I}
		\nocite{cameron2022-II}
	\end{frame}
	% Cover
	\begin{frame}
		\maketitle
	\end{frame}
\end{document}